\section{Background and Related Work}\label{sec:PreviousWork}
% Based on thesis sections 2 and 4.1--4.2. Claims checked against original sources.
% Source results and their implications for our design are distinguished below.
MCP separates the host application's coordination responsibilities from the capabilities exposed by servers~\cite{MCP2025Architecture}.
In this paper, dialogue state denotes the structured information retained across turns.
Four studies inform the use of tools, their orchestration and user control.

\textbf{ToolLLM.}
The framework addresses the training and evaluation of language models for real-world API use.
The ToolBench dataset covers 16,464 RESTful APIs, while a depth-first search method explores alternative sequences of calls.
The resulting ToolLLaMA model demonstrates multi-tool execution and generalisation to unseen APIs in the reported evaluation~\cite{qin2024toolllm}.
This supports studying tool selection and sequencing together.
Our prototype instead uses existing models with a small, predefined set of MCP tools and requires no model training.

\textbf{RestGPT.}
This approach separates task planning, API selection and execution.
The system uses API descriptions to construct calls and process responses, then feeds execution results back into planning.
On RestBench, covering TMDB and Spotify tasks, removing the planner reduces success rates~\cite{song2023restgpt}.
This motivates our comparison of orchestration with and without a separate planner.
The reported benefit applies to the evaluated benchmark and architecture, so its transfer to a bounded scheduling workflow requires evaluation.

\textbf{AI Chains.}
This approach connects multiple model calls so that each step can use earlier outputs.
An interactive interface exposes intermediate results and allows users to modify individual steps.
A study with 20 participants found improved task outcomes and perceived transparency, controllability and collaboration compared with a non-chaining interface~\cite{wu2022aichains}.
For our design, this motivates explicit intermediate state and opportunities for correction, while leaving the allocation of tool orchestration to be investigated.

\textbf{Guidelines for Human-AI Interaction.}
This work presents 18 design guidelines evaluated through several rounds, including a study with 49 practitioners examining 20 AI-enabled products.
The guidelines address communicating capabilities, correcting errors, clarifying uncertain goals and remembering recent interactions~\cite{amershi2019guidelines}.
These principles inform our dialogue design.
The requirement for explicit confirmation before booking is our application-specific implementation of human control.

Together, these studies connect tool capability, modular processing and human oversight.
Our study applies these perspectives to an MCP-based scheduling assistant and compares three implementations that distribute planning and execution between the LLM and backend.
Its focus is the resulting reliability and efficiency trade-offs within this controlled workflow.
